\documentclass[10pt,aspectratio=169]{beamer}
\usepackage{../../Common/course_slides}
\usepackage{../../Common/course_macros}
\usepackage[spanish,provide=*]{babel}

\title[Semana 2: Construcción del Arsenal Indicial]{De la Crisis Cartesiana al Arsenal de Índices}
\subtitle{Construyendo la notación de Einstein, $\delta_{ij}$, $\epsilon_{ijk}$ y Determinantes desde la Geometría}
\author{Geometría y Física I: Mecánica Clásica y Formas de Cartan}
\institute{}
\date{}

\begin{document}
\setlength{\abovedisplayskip}{3pt}
\setlength{\belowdisplayskip}{3pt}

% ---------------------------------------------------------------------
\begin{frame}
  \titlepage
\end{frame}

% ---------------------------------------------------------------------
\begin{frame}{Activación: Conexión con la Semana 01}
  \begin{block}{Recordando la Geometría de Vectores y Proyecciones}
    \footnotesize
    En la primera semana construimos las operaciones fundamentales en $\R^3$:
    \begin{itemize}
      \itemsep1pt
      \item Producto punto: $\vec{u}\cdot\vec{v} = u_1 v_1 + u_2 v_2 + u_3 v_3 = \|\vec{u}\|\|\vec{v}\|\cos\theta$.
      \item Descomposición ortogonal: $\vec{u}_\parallel = \left(\frac{\vec{u}\cdot\vec{v}}{\|\vec{v}\|^2}\right)\vec{v}$, \quad $\vec{u}_\perp = \vec{u} - \vec{u}_\parallel$.
      \item Producto cruz: $\vec{u}\times\vec{v} = (u_2 v_3 - u_3 v_2)\hat{e}_1 + (u_3 v_1 - u_1 v_3)\hat{e}_2 + (u_1 v_2 - u_2 v_1)\hat{e}_3$.
    \end{itemize}
  \end{block}
  \begin{alertblock}{Preguntas de Arranque para la Clase}
    \scriptsize
    \begin{enumerate}
      \itemsep1pt
      \item Al calcular productos triples como $\vec{A}\times(\vec{B}\times\vec{C})$, ¿cuántas multiplicaciones de componentes tenemos que escribir si usamos las fórmulas tradicionales?
      \item ¿Cómo podemos expresar la ortonormalidad de la base $\{\hat{e}_1, \hat{e}_2, \hat{e}_3\}$ de forma compacta y universal sin escribir matrices gigantes?
      \item ¿Es posible calcular identidades vectoriales complejas usando únicamente álgebra de índices y simetría?
    \end{enumerate}
  \end{alertblock}
\end{frame}

% ---------------------------------------------------------------------
\begin{frame}{Actividad 1: El Dolor Cartesiano y cómo Extraer Componentes}
  \begin{block}{Reto en clase (8 minutos)}
    \footnotesize
    Denotemos la base canónica como $\hat{e}_1, \hat{e}_2, \hat{e}_3$. Todo vector se expande como $\vec{A} = \sum_{j=1}^3 A_j\hat{e}_j$.
    \begin{enumerate}
      \itemsep1pt
      \item[(a)] Si calculas la 1\textsuperscript{a} componente de $\vec{A} \times (\vec{B} \times \vec{C})$ a mano:
      $
      [\vec{A} \times (\vec{B} \times \vec{C})]_1 = B_1(A_2 C_2 + A_3 C_3) - C_1(A_2 B_2 + A_3 B_3).
      $
      ¿Qué término debes sumar y restar artificialmente para completar $\vec{A}\cdot\vec{C}$ y $\vec{A}\cdot\vec{B}$?
      \item[(b)] Para extraer $A_i$ de $\vec{A}$ proyectando con $\hat{e}_i$: $\hat{e}_i \cdot \vec{A} = \sum_{j=1}^3 A_j (\hat{e}_i \cdot \hat{e}_j)$.
      ¿Cuánto vale $\hat{e}_i \cdot \hat{e}_j$ cuando $i = j$ y cuando $i \neq j$?
      \item[(c)] Usa tu respuesta de (b) para simplificar $\vec{A}\cdot\vec{B} = \left(\sum_i A_i\hat{e}_i\right)\cdot\left(\sum_j B_j\hat{e}_j\right)$.
    \end{enumerate}
  \end{block}
  \begin{alertblock}{Pregunta de Consolidación}
    \scriptsize
    ¿Por qué el producto escalar $\hat{e}_i \cdot \hat{e}_j$ actúa como un filtro que selecciona una única componente?
  \end{alertblock}
\end{frame}

% ---------------------------------------------------------------------
\begin{frame}{Discusión 1: Nacimiento de Einstein y de la Delta $\delta_{ij}$}
  \begin{columns}[c]
    \begin{column}{0.54\textwidth}
      \begin{block}{Construcción Geométrica de $\delta_{ij}$}
        \footnotesize
        No es una matriz arbitraria; es la \textbf{tabla de multiplicar de la base ortonormal}:
        \[
        \mathbf{\delta_{ij} \equiv \hat{e}_i \cdot \hat{e}_j} = \begin{cases} 1 & \text{si } i = j \\ 0 & \text{si } i \neq j \end{cases}
        \]
      \end{block}
      \footnotesize
      \textbf{Convenio de Suma de Einstein:} Todo índice repetido en un monomio se suma automáticamente de $1$ a $3$:
      \[
      A_i = \hat{e}_i \cdot (A_j\hat{e}_j) = A_j(\hat{e}_i\cdot\hat{e}_j) \implies \mathbf{\delta_{ij} A_j = A_i}
      \]
    \end{column}
    \begin{column}{0.43\textwidth}
      \begin{exampleblock}{Propiedades Fundamentales}
        \footnotesize
        \begin{itemize}
          \item \textbf{Producto punto:}
          \[
          \vec{A}\cdot\vec{B} = A_i B_j (\hat{e}_i\cdot\hat{e}_j) = \delta_{ij} A_i B_j = \mathbf{A_i B_i}
          \]
          \item \textbf{Traza (Dimensión del espacio $\R^3$):}
          \[
          \delta_{ii} = \delta_{11} + \delta_{22} + \delta_{33} = 1 + 1 + 1 = \mathbf{3}
          \]
          \item \textbf{Contracción de Deltas:}
          $\delta_{ik}\delta_{kj} = \delta_{ij}$.
        \end{itemize}
      \end{exampleblock}
    \end{column}
  \end{columns}
  \begin{exampleblock}{Pregunta de Consolidación}
    \scriptsize
    ¿Qué regla geométrica nos asegura que podemos renombrar cualquier índice mudo sin alterar la expresión física?
  \end{exampleblock}
\end{frame}

% ---------------------------------------------------------------------
\begin{frame}{Actividad 2: El Producto Cruz y el Símbolo de Levi-Civita}
  \begin{block}{Reto en clase (8 minutos)}
    \footnotesize
    Multipliquemos $\vec{A} = A_j\hat{e}_j$ y $\vec{B} = B_k\hat{e}_k$ usando únicamente la propiedad distributiva:
    \[
    \vec{A} \times \vec{B} = (A_j\hat{e}_j) \times (B_k\hat{e}_k) = A_j B_k (\hat{e}_j \times \hat{e}_k).
    \]
    \begin{enumerate}
      \item[(a)] Proyecta con $\hat{e}_i$ para extraer la $i$-ésima componente de $\vec{A}\times\vec{B}$:
      \[
      (\vec{A}\times\vec{B})_i = \hat{e}_i \cdot (\vec{A}\times\vec{B}) = A_j B_k \big[\hat{e}_i \cdot (\hat{e}_j \times \hat{e}_k)\big].
      \]
      \item[(b)] Definimos $\mathbf{\epsilon_{ijk} \equiv \hat{e}_i \cdot (\hat{e}_j \times \hat{e}_k)}$. Geométricamente, ¿qué representa el triple producto escalar de tres vectores?
      \item[(c)] Usando la interpretación de volumen, responde sin calcular: ¿Por qué $\epsilon_{112} = 0$? ¿Cuánto valen $\epsilon_{123}$ y $\epsilon_{132}$?
    \end{enumerate}
  \end{block}
  \begin{alertblock}{Pregunta de Consolidación}
    \scriptsize
    ¿Por qué cualquier intercambio de dos índices en $\epsilon_{ijk}$ invierte su signo algebraico?
  \end{alertblock}
\end{frame}

% ---------------------------------------------------------------------
\begin{frame}{Discusión 2: $\epsilon_{ijk}$ como Volumen Orientado de la Base}
  \begin{columns}[c]
    \begin{column}{0.56\textwidth}
      \begin{block}{Construcción Geométrica de $\epsilon_{ijk}$}
        \footnotesize
        $\mathbf{\epsilon_{ijk} \equiv \hat{e}_i \cdot (\hat{e}_j \times \hat{e}_k)} = \text{Vol}(\hat{e}_i, \hat{e}_j, \hat{e}_k)$
      \end{block}
      \footnotesize
      \begin{itemize}
        \itemsep0pt
        \item \textbf{Índices repetidos ($\epsilon_{112} = 0$):} Aristas paralelas $\implies$ \textbf{volumen nulo} ($21$ casos).
        \item \textbf{Permutación par ($(1,2,3), (2,3,1), (3,1,2)$):} Mano derecha $\implies \mathbf{+1}$ ($3$ casos).
        \item \textbf{Permutación impar ($(1,3,2), (2,1,3), (3,2,1)$):} Mano izquierda $\implies \mathbf{-1}$ ($3$ casos).
      \end{itemize}
      \begin{alertblock}{Producto Cruz Universal}
        \centering $\mathbf{(\vec{A} \times \vec{B})_i = \epsilon_{ijk} A_j B_k}$
      \end{alertblock}
    \end{column}
    \begin{column}{0.41\textwidth}
      \centering
      \begin{tikzpicture}[scale=0.72]
        \draw[line width=1.4pt, color=orangeacc] (0,0) circle (1.15cm);
        \node[circle, fill=blackmain, text=white, font=\bfseries, inner sep=2.5pt] at (90:1.15cm) {1};
        \node[circle, fill=blackmain, text=white, font=\bfseries, inner sep=2.5pt] at (330:1.15cm) {2};
        \node[circle, fill=blackmain, text=white, font=\bfseries, inner sep=2.5pt] at (210:1.15cm) {3};
        \draw[-{Stealth[length=2.2mm]}, line width=1.4pt, color=orangeacc] (68:1.45cm) arc (68:12:1.45cm) node[midway, right=1pt, font=\scriptsize\bfseries, color=orangeacc] {$+1$};
        \draw[-{Stealth[length=2.2mm]}, line width=1.4pt, color=orangeacc] (308:1.45cm) arc (308:252:1.45cm);
        \draw[-{Stealth[length=2.2mm]}, line width=1.4pt, color=orangeacc] (188:1.45cm) arc (188:132:1.45cm);
        \node[font=\scriptsize\bfseries, align=center] at (0,0) {Volumen\\Orientado};
      \end{tikzpicture}
    \end{column}
  \end{columns}
  \begin{exampleblock}{Pregunta de Consolidación}
    \scriptsize
    ¿Cómo se traduce la no-conmutatividad del producto cruz ($\vec{A}\times\vec{B} = -\vec{B}\times\vec{A}$) en la antisimetría $\epsilon_{ijk} = -\epsilon_{ikj}$?
  \end{exampleblock}
\end{frame}

% ---------------------------------------------------------------------
\begin{frame}{Actividad 3: La Identidad Maestra $\epsilon-\delta$}
  \begin{block}{Reto en clase (8 minutos)}
    \footnotesize
    Al calcular un doble producto cruz aparece la contracción $\epsilon_{ijk}\epsilon_{imn}$ sobre el índice común $i$.
    Fijemos $\mathbf{i = 1}$ y analicemos el término $\epsilon_{1jk}\epsilon_{1mn}$:
    \begin{enumerate}
      \item[(a)] Para que $\epsilon_{1jk} \neq 0$, los índices $(j, k)$ deben ser distintos de $1$ y no repetirse. ¿Cuáles son las dos únicas opciones ordenadas posibles para $(j, k)$?
      \item[(b)] Compara la pareja $(j, k)$ con $(m, n)$:
      \begin{itemize}
        \item Si coinciden en el mismo orden ($j = m$ y $k = n$), ¿cuánto vale el producto $\epsilon_{1jk}\epsilon_{1mn}$?
        \item Si están en orden invertido ($j = n$ y $k = m$), ¿cuánto vale el producto $\epsilon_{1jk}\epsilon_{1mn}$?
      \end{itemize}
      \item[(c)] Escribe una combinación de deltas de Kronecker que reproduzca exactamente $+1$ en el primer caso y $-1$ en el segundo.
    \end{enumerate}
  \end{block}
  \begin{alertblock}{Pregunta de Consolidación}
    \scriptsize
    ¿Por qué la contracción de dos símbolos antisimétricos da como resultado una diferencia de productos de deltas simétricas?
  \end{alertblock}
\end{frame}

% ---------------------------------------------------------------------
\begin{frame}{Discusión 3: Deducción y Estructura de $\epsilon_{ijk}\epsilon_{imn}$}
  \begin{columns}[c]
    \begin{column}{0.54\textwidth}
      \begin{block}{Identidad Maestra Universal $\epsilon-\delta$}
        \footnotesize
        \[
        \mathbf{\epsilon_{ijk}\epsilon_{imn} = \delta_{jm}\delta_{kn} - \delta_{jn}\delta_{km}}
        \]
      \end{block}
      \footnotesize
      \begin{itemize}
        \item \textbf{Término Positivo ($\delta_{jm}\delta_{kn}$):} Emparejamiento directo (primer índice libre con primero, segundo con segundo).
        \item \textbf{Término Negativo ($-\delta_{jn}\delta_{km}$):} Emparejamiento cruzado (invierte la orientación relativa).
      \end{itemize}
    \end{column}
    \begin{column}{0.43\textwidth}
      \begin{exampleblock}{Contracciones Adicionales}
        \footnotesize
        Haciendo $m = j$ en la Identidad Maestra:
        \[
        \epsilon_{ijk}\epsilon_{ijn} = \delta_{jj}\delta_{kn} - \delta_{jn}\delta_{kj} = 3\delta_{kn} - \delta_{kn} = \mathbf{2\delta_{kn}}
        \]
        Haciendo además $n = k$:
        \[
        \epsilon_{ijk}\epsilon_{ijk} = 2\delta_{kk} = 2(3) = \mathbf{6}
        \]
        (El número de permutaciones no nulas).
      \end{exampleblock}
    \end{column}
  \end{columns}
  \begin{exampleblock}{Pregunta de Consolidación}
    \scriptsize
    ¿Por qué $\epsilon_{ijk}\epsilon_{ijk} = 6$ coincide exactamente con el factorial $3!$?
  \end{exampleblock}
\end{frame}

% ---------------------------------------------------------------------
\begin{frame}{Actividad 4: La Gran Demolición de BAC--CAB y Potencia Magnética}
  \begin{block}{Reto en clase (8 minutos)}
    \footnotesize
    \begin{enumerate}
      \item[(a)] \textbf{Fórmula BAC--CAB:} Completa la deducción de la componente $i$ de $\vec{A} \times (\vec{B} \times \vec{C})$ usando la identidad maestra:
      \begin{align*}
        [\vec{A} \times (\vec{B} \times \vec{C})]_i &= \epsilon_{ijk} A_j (\vec{B}\times\vec{C})_k = \epsilon_{kij}\epsilon_{kmn} A_j B_m C_n \\
        &= (\delta_{im}\delta_{jn} - \delta_{in}\delta_{jm}) A_j B_m C_n = \dots
      \end{align*}
      \item[(b)] \textbf{Física de Lorentz:} La fuerza magnética sobre una carga es $F_i = q\,\epsilon_{ijk} v_j B_k$. La potencia instantánea es $P = \vec{F}\cdot\vec{v} = q\,\epsilon_{ijk} v_i v_j B_k$.
      ¿Qué ocurre con la contracción $\epsilon_{ijk} v_i v_j$ al contraer un tensor antisimétrico con uno simétrico?
    \end{enumerate}
  \end{block}
  \begin{alertblock}{Pregunta de Consolidación}
    \scriptsize
    ¿Cómo demostramos que una partícula cargada no cambia su energía cinética en presencia de un campo magnético estático sin resolver ecuaciones diferenciales?
  \end{alertblock}
\end{frame}

% ---------------------------------------------------------------------
\begin{frame}{Discusión 4: Solución, Simetría vs. Antisimetría y Síntesis}
  \begin{columns}[c]
    \begin{column}{0.52\textwidth}
      \begin{exampleblock}{(a) BAC--CAB Demostrada}
        \footnotesize
        \begin{align*}
          [\vec{A}\times(\vec{B}\times\vec{C})]_i &= (\delta_{im}\delta_{jn} - \delta_{in}\delta_{jm}) A_j B_m C_n \\
          &= B_i(A_j C_j) - C_i(A_j B_j) \\
          &= [\vec{B}(\vec{A}\cdot\vec{C}) - \vec{C}(\vec{A}\cdot\vec{B})]_i \quad \blacksquare
        \end{align*}
        ¡Obtenido directamente sin sumar ni restar términos arbitrarios a mano!
      \end{exampleblock}
    \end{column}
    \begin{column}{0.45\textwidth}
      \begin{block}{(b) Contracción Simétrica-Antisimétrica}
        \footnotesize
        Como $v_i v_j = v_j v_i$ (simétrico) y $\epsilon_{jik} = -\epsilon_{ijk}$ (antisimétrico):
        \[
        \mathbf{\epsilon_{ijk} v_i v_j = 0} \implies P = \vec{F}\cdot\vec{v} = 0.
        \]
        El campo magnético realiza trabajo nulo en todo instante.
      \end{block}
    \end{column}
  \end{columns}
  \begin{exampleblock}{Pregunta Final de Cierre}
    \scriptsize
    ¿Cuál es la ventaja fundamental de calcular con índices frente al álgebra vectorial tradicional de bachillerato?
  \end{exampleblock}
\end{frame}

\end{document}
